\documentclass{article}

\usepackage{../handouts}

\title{CSCI 535 Handout 05: Persistent Homology}
\date{10 November 2026}
\author{Name: \rule{3in}{0.15mm}}

\begin{document}
\maketitle

\section*{Warm-up with Homology ...}

As a class, do this warm-up set together.  This should take no more than 20
minutes. Those that feel comfortable with homology computation should take turns
presenting different parts.  Those who are learning this for the first time and
still gaining a footing, ask questions!

\begin{center}
\includegraphics[height=1.5in]{complex}
\end{center}

\begin{enumerate}
    \item Use $\Z_2$ coefficients and matrix reductions to compute the first
        homology group of the complex above.
        As a reminder, you'll need to
        take the following steps:
        \begin{enumerate}[itemsep=2in]
            \item Define $C_i$ for $i \in \N$ as being generated by the elements
                of $K_i$. (Note: using
                $\Z_2$ coefficients, we do not need to worry about picking an
                orientation; with other coefficient choices, we would need to
                pick an orientation for each simplex arbitrarily).  Fix an
                ordering of these generators.

            \pagebreak
            \item Give the boundary matrices for the linear operators
                $\partial_1$, $\partial_2$.  When you write them out, leave
                plenty of space for reducing this later!
                \vspace{4in}
            \item Reduce only using column-wise additions for now. (Either in
                place above or step it out for each column operation).
                The reduced labels the zero columns give the generators of
                $Z_1$ and $Z_2$.  What are they (if none, that means the group
                is the trivial group)?

            \pagebreak
            \item Continue the reduction by doing row-wise additions.  The
                labels of the non-zero rows are the generators of $B_0$ and
                $B_1$.  What are they?

            \item Note that we also know $Z_0$ and $B_2$ are the trivial groups.
                Knowing that $\Z_2^i / Z_2^j \cong \Z_2^{i-j}$, what are the
                homology groups, up to isomorphism? (i.e., I don't need the exact
                generators, just something along the lines of: $H_{30} \cong
                \Z_2^4$).

            \item HW Challenge: The homology groups you get do not change if you do the
                operations in different orders.  But, you might get different
                generators.  Try to do different column operations to get
                different generators of $Z_1$. Explain why it is still the same
                $Z_1$.
                \vspace{2in}
        \end{enumerate}
\end{enumerate}

\pagebreak
\section*{Homology and Filtrations}

This section of the class will be the majority of the class period. Reserve
about 10 minutes at the end to come back as a class.

Colloquially, we describe homology as \emph{connected components}
(zero-dimensional), \emph{tunnels} or \emph{loops} (one-dimensional), and
\emph{voids} (two-dimensional).  This is because $H_0$ captures connected
components.  In~$\R^2$,~$H_1$ captures loops; in $\R^3$, it captures tunnels.
In $\R^3$, $H_2$ captures voids, or `insides' (e.g., the torus creates one
`inside' a torus and a sphere create two `insides').  With that in mind, let's
see what happens to homology as we follow a couple filtrations.

Note: if you notice that if $K$ is a geometric simplicial complex in $\R^3$,
then the voids are the number of (path) connected components of $\R^3 \setminus
K$, then you are observing duality! Cool beans.

We've looked at \v{C}ech complexes complexes before, and how they
form a nested sequence of simplicial complexes as the radius increases (same for
Rips complexes!) In general, a \emph{filtration} is a nested sequence of
simplicial complexes. A \emph{filter function} $f \colon K \to \R$ assigns a
function value to each simplex such that $F_t:=f^{-1}(-\infty,t]$ is always a
simplicial complex (then, $\{F_t\}_t\in \R$ is the corresponding function).  We
use $t$ as the parameter here, as we often think of it as \emph{time}.

For this part of the class, consider the following filter functions built on the
same underlying complex, the complete complex on four vertices
($K_0=\{a,b,c\})$:
\begin{enumerate}
    \item[Group 1:]
        $f(a)=1$,
        $f(b)=2$,
        $f(c)=3$,
        $f(ac)=4$,
        $f(ab)=5$,
        $f(bc)=6$, and
        $f(abc)=7$.
    \item[Group 2:]
        $f(a)=0$,
        $f(b)=0$,
        $f(c)=0$,
        $f(ac)=\frac{\sqrt{3.25}}{2} \approx 0.9$,
        $f(ab)=\frac{\sqrt{3.25}}{2}$,
        $f(bc)=1$, and
        $f(abc) =1$.
    \item[Group 3:]
        $f(a)=0$,
        $f(b)=0$,
        $f(c)=0$,
        $f(ac)=\frac{\sqrt{3.25}}{2}$,
        $f(ab)=\frac{\sqrt{3.25}}{2}$,
        $f(bc)=1$, and
        $f(abc)=\frac{13}{12}\approx 1.08 $.
    \item[Bonus:]
        $f(a)=-1$,
        $f(b)=0$,
        $f(c)=1$,
        $f(ab)=0$,
        $f(bc)=1$,
        $f(ac)=1$,
        $f(abc)=1$.
\end{enumerate}
Note: here, we are using the convention that $ab$ is the simplex containing the
vertices $a$ and $b$, etc.


\begin{enumerate}[itemsep=2in]
    \item Draw your filtration (luckily, there are only a finite number of
        complexes that arise, even though every real number gives you a
        simplicial complex, so we can draw all steps as follows: create a
        `number line' and mark off each time the complex changes. Then, draw the complex at
        each `changepoint').
    \item Name your filtration. (Could it be a \v{C}ech or a Rips complex? If
        neither, how would you describe what is happening in it?)
    \item For each complex in the filtration, state what the homology class is.
        Give a (minimal) generating set for the homology class (minimal here =
        using the least number of generators).
    \item At what times is a new
        homology generator (or multiple new ones) introduced?
        These are our homology `births'.
    \item At what times do two (or more) connected components merge together?
        These are the zero-dimensional deaths.
    \item At what times do two (or more) non-homologous zero-cycles merge / become homologous?
        (Note: this could happen as something nontrivial becoming trivial /
        `filling in', or as
        two nontrivial things becoming homologous).
        These are the one-dimensional deaths.
    \item Persistence pairs the births and deaths using the `eldest survives'
        rule: if two classes merge, the label of the elder class `lives on' and
        the younger class is the one that `died'. Ties are broken
        arbitrarily. Using this rule, what are your persistence pairs? (Note:
        there should be three pairs, and one unpaired vertex).
    \item Fun fact: persistence computation is pretty much the same as homology
        (well, complexity-wise), we just have to be careful about what order we
        do our matrix operations (so it makes sense with respect to the `time').
        In addition, if we don't
        swap columns/rows, then we can associate `times' with the columns and
        rows. Process the
        columns left to right. If the lowest non-zero entry of a column is
        `unclaimed', then claim it! You've found a persistence pairing (the
        times corresponding to the given column and row). If the lowest non-zero
        entry of the column was already claimed, add the column that claimed it
        to the current column.  Now, the lowest 1 is necessarily `above' where
        we were.  Repeat until the column can claim something or becomes zero.
        (Hey! can you prove this terminates using a decrementing function?). If
        time allows, try
        finding the pairings you found above using matrix reductions.
\end{enumerate}

\pagebreak
\section*{}

In the last 10-ish minutes of class, discuss the following:

\begin{enumerate}
    \item What is the relationship between Group 2 and Group 3's filtrations?
        Is it possible to think of $a,b,c$ as points in the plane? If so, what
        points?
    \item What is the relationship between Group 1 and Group 3's filtrations?
        (Hint: draw them on a `number line' on top of each other, then draw a
        line from each complex in Group 1's filtration to the first time that complex
        is seen in Group 3's filtration, as a subcomplex).
    \item Give an example of two nonhomologous, nontrivial one-cycles becoming
        homologous in a filtration. (i.e., two things that did not differ by a
        bounding cycle before time $t$, now, differ by a bounding cycle after
        time $t$).
\end{enumerate}

\end{document}
